% from ega4-20-fr.tex
\oldpage[IV-4]{254}
  en effet trivialement nécessaire.
  Inversement, si elle est vérifiée, et si $x$ est un point quelconque de $X$, $s$ son image dans $S$, il existe par hypothèse un voisinage ouvert $U$ de $x$ dans $X$ et deux sections $u$, $t$ de $\sh{O}_X$ au-dessus de $U$ telles que $t\in T_{X/S}(U)$ et $\varphi|U=u/t$ ; l'hypothèse entraîne que si $u_s$ est l'image de $u$ dans $\Gamma(U\cap X_s,\sh{O}_{X_s})$, $u_s$ est régulière au point $x$.
  En restreignant $U$, on peut donc supposer, en vertu de \sref{IV.11.3.8}, que $u\in T_{X/S}(U)$, d'où la conclusion.

  Lorsque $\varphi$ est inversible dans $M(X/S)$, on dit encore que $\varphi$ est une \emph{fonction méromorphe régulière relative à $S$}.
  On notera que $\varphi\in M(X/S)$ peut être inversible dans $M(X)$ (autrement dit, être \emph{régulière} au sens de \sref{IV.20.1.8}) sans l'être dans $M(X/S)$, comme le montre aussitôt l'exemple de~(i).

  \item Soit $\sh{L}$ un $\sh{O}_X$-Module \emph{inversible}, et soit $\varphi$ une section méromorphe \emph{régulière} de $\sh{L}$ au-dessus de $X$ \sref{IV.20.1.8} ; on dit que $\varphi$ est \emph{régulière relativement à $S$} si, pour tout ouvert $U$ de $X$ tel que $\sh{L}|U$ soit isomorphe à $\sh{O}_U$, $\varphi|U$ correspond à un élément de $\Gamma(U,\sh{M}_X)$ qui est \emph{régulière relativement à $S$} ; en vertu de~(ii), il est immédiat qu'il faut et il suffit pour cela que, pour tout $s\in S$, $\varphi_s$ soit une section méromorphe \emph{régulière} \sref{IV.20.1.8} du $\sh{O}_{X_s}$-Module inversible $\sh{L}_s=\sh{L}\otimes_{\sh{O}_X}k(s)$.
  Si $\varphi'$ est l'inverse de $\varphi$ dans $\sh{L}^{-1}$ \sref{IV.20.1.10}, $\varphi'$ est alors aussi régulière relative à $S$.
  Si $\sh{L}_1$ est un second $\sh{O}_X$-Module inversible, $\varphi_1$ une section méromorphe de $\sh{L}_1$ au-dessus de $X$, régulière relativement à $S$, alors $\varphi\otimes\varphi_1$ est une section méromorphe de $\sh{L}\otimes\sh{L}_1$ au-dessus de $X$, régulière relativement à $S$.
\end{enumerate}
\end{remarks}

\begin{proposition}[20.6.6]
\label{IV.20.6.6}
Soient $X$ un $S$-préschéma plat et localement de présentation finie sur $S$, $\sh{F}$ un $\sh{O}_X$-Module localement libre de type fini ; pour tout $s\in S$, on note $X_s$ la fibre au point $s$ du morphisme structural $f:X\to S$.
Soit $\varphi$ une section méromorphe de $\sh{F}$ au-dessus de $X$, relativement à $S$, et supposons que $\varphi$ soit définie en tous les points $x\in X$ tels que $\operatorname{prof}(\sh{O}_{X_{f(x)},x})=1$.
Alors $\varphi$ est partout définie.
\end{proposition}

\begin{proof}
Par hypothèse, $\operatorname{dom}(\varphi)\cap X_s$ est schématiquement dense dans $X_s$ pour tout $s\in S$, donc contient les points $x$ de $X_s$ tels que $\operatorname{prof}(\sh{O}_{X_s,x})=0$ \sref{IV.5.10.2} ; l'hypothèse signifie donc que si l'on pose $Z=X-\operatorname{dom}(\varphi)$, on a $\operatorname{prof}(\sh{O}_{X_{f(x)},x})\geq2$ en tous les points de $Z$.
Il suffit donc d'appliquer \sref{IV.19.9.8}.
\end{proof}

\begin{env}[20.6.7]
\EGAReaderTarget{ega4.c04.statement.20-6-7}{IV.20.6.7}
Soient $X$, $X'$ deux $S$-préschémas plats et localement de présentation finie sur $S$, $f=(\psi,\theta):X'\to X$ un $S$-morphisme.
Pour tout ouvert $U$ de $X$, notons $T_f(U)$ l'ensemble des sections $t\in T_{X/S}(U)$ dont l'image dans $\Gamma(f^{-1}(U),\sh{O}_{X'})$ appartienne à $T_{X'/S}(f^{-1}(U))$ ; il est immédiat que $U\mapsto T_f(U)$ est un sous-faisceau du faisceau d'ensembles $\sh{T}_{X/S}$, que l'on note $\sh{T}_f$.
On pose $\sh{M}_{X/S,f}=\sh{O}_X[\sh{T}_f^{-1}]$ ; c'est un sous-faisceau d'anneaux de $\sh{M}_{X/S}$, et on en déduit canoniquement de $\theta^\#: \psi^*(\sh{O}_X)\to\sh{O}_{X'}$ un homomorphisme de faisceaux d'anneaux $\theta'^\#: \psi^*(\sh{M}_{X/S,f})\to\sh{M}_{X'/S}$ prolongeant $\theta^\#$.
Si une fonction méromorphe $\varphi$ sur $X$, relative à $S$, est \emph{une section de $\sh{M}_{X/S,f}$}, $\Gamma(\theta'^\#)(\varphi)$ est une fonction méromorphe sur $X'$, dite \emph{image réciproque de $\varphi$ par $f$}, et notée $\varphi\circ f$ si cela n'entraîne pas confusion.
On étend aussitôt de même les définitions de \sref{IV.20.1.11} relatives aux $\sh{O}_X$-Modules.
\end{env}

\begin{proposition}[20.6.8]
\label{IV.20.6.8}
Avec les notations de \sref{IV.20.6.7}, si le $S$-morphisme $f:X'\to X$ est \emph{plat}, on a $\sh{M}_{X/S,f}=\sh{M}_{X/S}$, et l'homomorphisme $\varphi\mapsto\varphi\circ f$ est défini dans $M(X/S)$ tout entier.
\end{proposition}

\begin{proof}
